\section{Threat Model and Measurement Protocol}
\label{sec:protocol}

\subsection{Adversary}

The attacker controls the message text and nothing else: no access to the model
weights, gradients, logits or detector code, and no way to retrain it. The
attacker also wants the message to stay readable, so the edit budget is capped:
at most 15\% of the message's units (characters for the character-level attack,
whitespace-separated tokens for the word-level attack) may change, and the change
is measured, not estimated.

\paragraph{Two channels, reported separately.} Not every attack in this paper
assumes the same interface, and the paper never adds them together. The
\emph{verdict channel} returns the one-word generated answer: malicious or
benign, one bit per query. That is what a spammer probing a filter receives, and
it is the information source for the character and random-word attacks
(Section~\ref{sec:results}). The \emph{probability channel} also returns
$P(\text{malicious})$ --- here by scoring the two permitted answer tokens --- and
only the saliency-guided attack uses it (Section~\ref{sec:saliency}), because a ranking
needs a number and one bit per candidate cannot give one. Every row in every
table records which channel produced it, the two channels are never mixed inside
one comparison, and the disagreement between them is reported on its own
(Section~\ref{sec:discussion}). The ablation of Section~\ref{sec:ablation} holds the
information source fixed across arms, precisely so that the cost of dropping from
the probability channel to the verdict channel is measured rather than assumed.

The verdict-only adversary is deliberately weak, and for that reason its ASR is
a lower bound. An attacker with weight access, an
embedding-space attack or a saliency-guided word search would do better. A weak
adversary is useful precisely because a flip under it is unambiguously bad news
for the detector: nothing exotic happened, the message simply stopped looking
suspicious.

\subsection{The four rules that make an ASR quotable}

\paragraph{Rule 1 --- count only the messages the model had caught, and check them
again.}
Only messages the clean model classified \emph{malicious} are attacked. The
original text is queried again inside the attack run instead of trusting the
baseline label, and a sample whose original no longer reads as malicious is
\emph{excluded} from the total rather than counted as a failed attack. In our
runs this happened 52 times, and each exclusion is recorded with its reason.

\paragraph{Where the 52 exclusions come from.} They fall into exactly two kinds.
Forty-three are samples that the saliency attack's probability channel read as
benign on the original --- all of them on Qwen2.5-3B, a channel problem discussed
in Section~\ref{sec:discussion}. The other nine are word-attack samples whose edited
text came out identical to the original, so no edit was made and there was
nothing to score. No original was dropped for any other reason: the verdict
channel never re-read a baseline-caught message as benign, so the character and
saliency runs on the other two models exclude nothing at all.

\paragraph{Rule 2 --- the same prompt on both sides.}
The baseline and the attack call the identical message template, including the
sample's own task description, so an email is never scored with the SMS prompt.
A baseline-versus-attack comparison across two different prompts would measure
the prompt, not the attack.

\paragraph{Rule 3 --- one shared budget, measured exactly.}
The 15\% budget is a single allowance for the whole edit, not one per operation,
and the edit rate is the exact Levenshtein distance (the smallest number of
single-character insertions, deletions and substitutions that turns one string
into the other) divided by the length of the original. Both details matter.
Giving each operation its own allowance let the real rate reach 22\% while the
reported budget was 15\%; and estimating the distance from a `difflib' opcode tally
reported 204 edits for 56 deleted characters on a repetitive message, because
`difflib' does not compute a minimal edit script.

\paragraph{Rule 4 --- name the similarity you actually computed.}
We report word-set Jaccard overlap and call it Jaccard. No sentence-similarity
model is used anywhere in this study, so no BERTScore is claimed. For the
character attack we add that Jaccard is the wrong tool: a look-alike substitution
changes the token, so word-set overlap collapses even though a human reads the
same message.

\subsection{Perturbations: what each attack changes}

\textbf{Character level.} Four operations share the budget: look-alike
substitution (the Latin letters `a', `e', `o', `p', `c', `x', `y', `i', `s'
replaced by Cyrillic characters that print the same), random deletion, random
insertion of lowercase letters, and swapping two neighbouring characters
(counting as two edits).

\textbf{Word level.} Four operations share the budget: replacing a word with one
of its listed synonyms from a hand-written 14-entry table, word deletion, word
duplication, and word swap. The synonym table is deliberately small, which
weakens the attack and strengthens the negative result.

\textbf{How many attempts, and what that costs.} The two random attacks give each
sample up to three independent attempts from a per-sample seeded random
generator, stopping at the first flip, so one message can cost one to four
queries: one for the re-queried original plus one per attempt. On the recorded
runs that comes to 3.9 queries per message and 8{,}047 queries for the 2{,}079
attacked messages --- the number of messages and the number of queries are
different quantities and are reported separately throughout. The saliency attack
is not attempt-based: it probes every candidate position once and then re-queries
after each accepted edit, 22.9 queries per message on average. The attempt column
in every row carries how many tries a message actually took, which is why the
count above is measured rather than assumed.

Every attempt's edited text, verdict, edit count and rate are written to the
per-sample CSV, so the tables can be rebuilt without the model.

\subsection{Statistics}

An ASR of $k$ flips out of $n$ attacks is a share of a fixed number of attempts.
We report the 95\% Wilson score interval, which stays inside $[0,1]$ and, more
importantly, does not collapse at the edge: a row with $k=0$ of $n=188$ has upper
bound 2.0\%, not 0\%. As a cross-check we also compute a percentile bootstrap
over the same rows (10{,}000 resamples, fixed seed), which agrees with Wilson
wherever $k>0$; where $k=0$ the bootstrap collapses to $[0,0]$ and we quote
Wilson instead. The interval covers sampling variation over \emph{messages}, not
variation of the perturbation seed; measuring the latter needs a second run over
the same samples, which is what the seed sweep below does.

\paragraph{Seed variation.} The intervals above cover variation over messages,
not over the perturbation generator. The runner therefore takes \texttt{--seed}
and a sweep driver (\texttt{src/run\_seed\_sweep.py}) re-runs one arm at several
seeds, so the spread between seeds can be reported next to the sampling interval
rather than absorbed into it. The sweep is implemented; its cells are not part of
the numbers in this version, and Section~\ref{sec:discussion} lists that as a gap
rather than implying the intervals cover it.
